\begin{register}{H}{EFUSE\_PGM\_DATA\regindex{n}\_REG (\regindex{n}: 0-7)}{0x0000+0x4*\regindex{n}}\label{regdesc:EFUSEPGMDATANREG}
    \regfield{EFUSE\_PGM\_DATA\_0}{32}{0}{{0x000000}}%
    \reglabel{Reset}\regnewline%
    \begin{regdesc}\begin{reglist}
    \label{fielddesc:EFUSEPGMDATAN}\item [EFUSE\_PGM\_DATA\_\regindex{n}] 配置待烧写数据的第 \regindex{n} 个 32 位数据。(R/W)
    \end{reglist}\end{regdesc}
    \end{register}


    \begin{register}{H}{EFUSE\_PGM\_CHECK\_VALUE\regindex{n}\_REG (\regindex{n}: 0-2)}{0x0020+0x4*\regindex{n}}\label{regdesc:EFUSEPGMCHECKVALUENREG}
    \regfield{EFUSE\_PGM\_RS\_DATA\_0}{32}{0}{{0x000000}}%
    \reglabel{Reset}\regnewline%
    \begin{regdesc}\begin{reglist}
    \label{fielddesc:EFUSEPGMRSDATAN}\item [EFUSE\_PGM\_RS\_DATA\_\regindex{n}] 配置待烧写的第 \regindex{n} 个 32 位 RS 码。(R/W)
    \end{reglist}\end{regdesc}
    \end{register}


    \begin{register}{H}{EFUSE\_RD\_WR\_DIS0\_REG}{0x002C}\label{regdesc:EFUSERDWRDIS0REG}
    \regfield{EFUSE\_WR\_DIS}{32}{0}{{0x000000}}%
    \reglabel{Reset}\regnewline%
    \begin{regdesc}\begin{reglist}
    \label{fielddesc:EFUSEWRDIS}\item [EFUSE\_WR\_DIS] 表示是否使能 eFuse 存储器各个位的烧写。\\
    1：不使能\\
    0：使能\\
    (RO)
    \end{reglist}\end{regdesc}
    \end{register}


    \begin{register}{H}{EFUSE\_RD\_REPEAT\_DATA0\_REG}{0x0030}\label{regdesc:EFUSERDREPEATDATA0REG}
    \regfield{(reserved)}{3}{29}{000}%
    \regfield{EFUSE\_SECURE\_BOOT\_KEY\_REVOKE\_2}{1}{28}{{0}}%
    \regfield{EFUSE\_SECURE\_BOOT\_KEY\_REVOKE\_1}{1}{27}{{0}}%
    \regfield{EFUSE\_SECURE\_BOOT\_KEY\_REVOKE\_0}{1}{26}{{0}}%
    \regfield{EFUSE\_SPI\_BOOT\_CRYPT\_CNT}{3}{23}{{0x0}}%
    \regfield{EFUSE\_WDT\_DELAY\_SEL}{2}{21}{{0x0}}%
    \regfield{EFUSE\_VDD\_SPI\_AS\_GPIO}{1}{20}{{0}}%
    \regfield{EFUSE\_USB\_EXCHG\_PINS}{1}{19}{{0}}%
    \regfield{EFUSE\_USB\_DREFL}{2}{17}{{0x0}}%
    \regfield{EFUSE\_USB\_DREFH}{2}{15}{{0x0}}%
    \regfield{EFUSE\_DIS\_DOWNLOAD\_MANUAL\_ENCRYPT}{1}{14}{{0}}%
    \regfield{EFUSE\_DIS\_PAD\_JTAG}{1}{13}{{0}}%
    \regfield{EFUSE\_JTAG\_SEL\_ENABLE}{1}{12}{{0}}%
    \regfield{EFUSE\_SPI\_DOWNLOAD\_MSPI\_DIS}{1}{11}{{0}}%
    \regfield{EFUSE\_DIS\_FORCE\_DOWNLOAD}{1}{10}{{0}}%
    \regfield{EFUSE\_DIS\_USB\_SERIAL\_JTAG}{1}{9}{{0}}%
    \regfield{EFUSE\_DIS\_USB\_JTAG}{1}{8}{{0}}%
    \regfield{EFUSE\_DIS\_ICACHE}{1}{7}{{0}}%
    \regfield{EFUSE\_RD\_DIS}{7}{0}{{0x0}}%
    \reglabel{Reset}\regnewline%
    \begin{regdesc}\begin{reglist}
    \label{fielddesc:EFUSERDDIS}\item [EFUSE\_RD\_DIS] 表示是否使能 eFuse 各个块 (BLOCK4 \verb+~+ BLOCK10) 的读取。\\
    1：不使能\\
    0：使能\\
    (RO)
    \label{fielddesc:EFUSEDISICACHE}\item [EFUSE\_DIS\_ICACHE] 表示是否使能 iCache。\\
    1：不使能\\
    0：使能\\
    (RO)
    \label{fielddesc:EFUSEDISUSBJTAG}\item [EFUSE\_DIS\_USB\_JTAG] 表示是否使能 USB 转 JTAG 功能。\\
    1：不使能\\
    0：使能\\
    (RO)
    \label{fielddesc:EFUSEDISUSBSERIALJTAG}\item [EFUSE\_DIS\_USB\_SERIAL\_JTAG] 表示是否使能 USB-Serial-JTAG。\\
    1：不使能\\
    0：使能\\
    (RO)
    \label{fielddesc:EFUSEDISFORCEDOWNLOAD}\item [EFUSE\_DIS\_FORCE\_DOWNLOAD] 表示是否使能强制芯片进入 Download 模式。\\
    1：不使能\\
    0：使能\\
    (RO)

    \item [见下页……]
    \end{reglist}\end{regdesc}
    \end{register}
    \addtocounter{Regfloat}{-1}
    \begin{register}{H}{EFUSE\_RD\_REPEAT\_DATA0\_REG}{0x0030}
    \begin{regdesc}\begin{reglist}
    \item [接上页……]

    \label{fielddesc:EFUSESPIDOWNLOADMSPIDIS}\item [EFUSE\_SPI\_DOWNLOAD\_MSPI\_DIS] 表示 boot\_mode\_download 过程中是否使能 SPI0 控制器。\\
    1：不使能\\
    0：使能\\
    (RO)

    \label{fielddesc:EFUSEJTAGSELENABLE}\item [EFUSE\_JTAG\_SEL\_ENABLE] 表示 \hyperref[fielddesc:EFUSEDISPADJTAG]{EFUSE\_DIS\_PAD\_JTAG} 和 \hyperref[fielddesc:EFUSEDISUSBJTAG]{EFUSE\_DIS\_USB\_JTAG} 均配置为 0 时，是否使能通过 GPIO15 的 strapping 值选择 JTAG 信号源 usb\_to\_jtag 与 pad\_to\_jtag。\\
    1：使能\\
    0：不使能\\
    (RO)
    \label{fielddesc:EFUSEDISPADJTAG}\item [EFUSE\_DIS\_PAD\_JTAG] 表示是否硬禁用 JTAG（永久）。\\
    1：禁用\\
    0：不禁用\\
    (RO)
    \label{fielddesc:EFUSEDISDOWNLOADMANUALENCRYPT}\item [EFUSE\_DIS\_DOWNLOAD\_MANUAL\_ENCRYPT] 表示是否使能 flash 加密（SPI 启动模式除外）。\\
    1：禁用\\
    0：不禁用\\
    (RO)
    \label{fielddesc:EFUSEUSBDREFH}\item [EFUSE\_USB\_DREFH] 表示单端输入阈值 Vrefh，范围为 1.76 V 到 2 V，步长为 80 mV。(RO)
    \label{fielddesc:EFUSEUSBDREFL}\item [EFUSE\_USB\_DREFL] 表示单端输入阈值 Vrefl，范围为 1.76 V 到 2 V，步长为 80 mV。(RO)
    \label{fielddesc:EFUSEUSBEXCHGPINS}\item [EFUSE\_USB\_EXCHG\_PINS] 表示是否交换 D+ 和 D- 管脚。\\
    1：交换\\
    0：不交换\\
    (RO)
    \label{fielddesc:EFUSEVDDSPIASGPIO}\item [EFUSE\_VDD\_SPI\_AS\_GPIO] 表示是否将 VDD\_SPI 管脚用作普通 GPIO。\\
    1：用作普通 GPIO\\
    0：不用作普通 GPIO\\
    (RO)
    \label{fielddesc:EFUSEWDTDELAYSEL}\item [EFUSE\_WDT\_DELAY\_SEL] 表示 RTC 看门狗 STG0 超时的阈值等级。\\
    0：原始阈值的 2 倍 \\
    1：原始阈值的 4 倍 \\
    2：原始阈值的 8 倍 \\
    3：原始阈值的 16 倍 \\
    (RO)
    \label{fielddesc:EFUSESPIBOOTCRYPTCNT}\item [EFUSE\_SPI\_BOOT\_CRYPT\_CNT] 表示是否使能 SPI boot 加解密。\\
    奇数个比特为 1：使能\\
    偶数个比特为 1：不使能\\
    (RO)

    \item [见下页……]
    \end{reglist}\end{regdesc}
    \end{register}

    \addtocounter{Regfloat}{-1}
    \begin{register}{H}{EFUSE\_RD\_REPEAT\_DATA0\_REG}{0x0030}
    \begin{regdesc}\begin{reglist}
    \item [接上页……]

    \label{fielddesc:EFUSESECUREBOOTKEYREVOKE0}\item [EFUSE\_SECURE\_BOOT\_KEY\_REVOKE\_0] 表示是否撤销第一个安全启动秘钥。\\
    1：撤销\\
    0：不撤销\\
    (RO)
    \label{fielddesc:EFUSESECUREBOOTKEYREVOKE1}\item [EFUSE\_SECURE\_BOOT\_KEY\_REVOKE\_1] 表示是否撤销第二个安全启动秘钥。\\
    1：撤销\\
    0：不撤销\\
    (RO)
    \label{fielddesc:EFUSESECUREBOOTKEYREVOKE2}\item [EFUSE\_SECURE\_BOOT\_KEY\_REVOKE\_2] 表示是否撤销第三个安全启动秘钥。\\
    1：撤销\\
    0：不撤销\\
    (RO)
    \end{reglist}\end{regdesc}
    \end{register}


    \begin{register}{H}{EFUSE\_RD\_REPEAT\_DATA1\_REG}{0x0034}\label{regdesc:EFUSERDREPEATDATA1REG}
    \regfield{EFUSE\_FLASH\_TPUW}{4}{28}{{0x0}}%
    \regfield{EFUSE\_SECURE\_BOOT\_AGGRESSIVE\_REVOKE}{1}{27}{{0}}%
    \regfield{EFUSE\_SECURE\_BOOT\_EN}{1}{26}{{0}}%
    \regfield{EFUSE\_SEC\_DPA\_LEVEL}{2}{24}{{0x0}}%
    \regfield{EFUSE\_KEY\_PURPOSE\_5}{4}{20}{{0x0}}%
    \regfield{EFUSE\_KEY\_PURPOSE\_4}{4}{16}{{0x0}}%
    \regfield{EFUSE\_KEY\_PURPOSE\_3}{4}{12}{{0x0}}%
    \regfield{EFUSE\_KEY\_PURPOSE\_2}{4}{8}{{0x0}}%
    \regfield{EFUSE\_KEY\_PURPOSE\_1}{4}{4}{{0x0}}%
    \regfield{EFUSE\_KEY\_PURPOSE\_0}{4}{0}{{0x0}}%
    \reglabel{Reset}\regnewline%
    \begin{regdesc}\begin{reglist}
    \label{fielddesc:EFUSEKEYPURPOSE0}\item [EFUSE\_KEY\_PURPOSE\_0] 表示 Key0 用途。 (RO)
    \label{fielddesc:EFUSEKEYPURPOSE1}\item [EFUSE\_KEY\_PURPOSE\_1] 表示 Key1 用途。 (RO)
    \label{fielddesc:EFUSEKEYPURPOSE2}\item [EFUSE\_KEY\_PURPOSE\_2] 表示 Key2 用途。 (RO)
    \label{fielddesc:EFUSEKEYPURPOSE3}\item [EFUSE\_KEY\_PURPOSE\_3] 表示 Key3 用途。 (RO)
    \label{fielddesc:EFUSEKEYPURPOSE4}\item [EFUSE\_KEY\_PURPOSE\_4] 表示 Key4 用途。 (RO)
    \label{fielddesc:EFUSEKEYPURPOSE5}\item [EFUSE\_KEY\_PURPOSE\_5] 表示 Key5 用途。 (RO)
    \label{fielddesc:EFUSESECDPALEVEL}\item [EFUSE\_SEC\_DPA\_LEVEL] 表示是否通过配置时钟随机分频模式来确定 DPA 安全等级。\\
    0：没有作用\\
    1：通过\\
    (RO)
    \label{fielddesc:EFUSESECUREBOOTEN}\item [EFUSE\_SECURE\_BOOT\_EN] 表示是否使能安全启动。\\
    1：使能\\
    0：不使能\\
    (RO)
    \label{fielddesc:EFUSESECUREBOOTAGGRESSIVEREVOKE}\item [EFUSE\_SECURE\_BOOT\_AGGRESSIVE\_REVOKE] 表示是否使能安全启动的激进撤销。\\
    1：使能\\
    0：不使能\\
    (RO)
    \label{fielddesc:EFUSEFLASHTPUW}\item [EFUSE\_FLASH\_TPUW] 表示 flash 上电等待时间。单位：ms。当值小于 15 时，等待时间为配置值。当值大于等于 15 时，等待时间为 2 倍配置值。(RO)
    \end{reglist}\end{regdesc}
    \end{register}


    \begin{register}{H}{EFUSE\_RD\_REPEAT\_DATA2\_REG}{0x0038}\label{regdesc:EFUSERDREPEATDATA2REG}
    \regfield{EFUSE\_ECC\_FORCE\_CONST\_TIME}{1}{31}{{0}}%
    \regfield{EFUSE\_ECDSA\_DISABLE\_P192}{1}{30}{{0}}%
    \regfield{EFUSE\_DIS\_WIFI6}{1}{29}{{0}}%
    \regfield{EFUSE\_XTS\_DPA\_PSEUDO\_LEVEL}{2}{27}{{0x0}}%
    \regfield{EFUSE\_XTS\_DPA\_CLK\_ENABLE}{1}{26}{{0}}%
    \regfield{EFUSE\_HYS\_EN\_PAD}{1}{25}{{0}}%
    \regfield{EFUSE\_SECURE\_BOOT\_DISABLE\_FAST\_WAKE}{1}{24}{{0}}%
    \regfield{EFUSE\_SECURE\_VERSION}{16}{8}{{0x00}}%
    \regfield{EFUSE\_FORCE\_SEND\_RESUME}{1}{7}{{0}}%
    \regfield{EFUSE\_UART\_PRINT\_CONTROL}{2}{5}{{0x0}}%
    \regfield{EFUSE\_ENABLE\_SECURITY\_DOWNLOAD}{1}{4}{{0}}%
    \regfield{EFUSE\_DIS\_USB\_SERIAL\_JTAG\_DOWNLOAD\_MODE}{1}{3}{{0}}%
    \regfield{EFUSE\_DIS\_USB\_SERIAL\_JTAG\_ROM\_PRINT}{1}{2}{{0}}%
    \regfield{EFUSE\_DIS\_DIRECT\_BOOT}{1}{1}{{0}}%
    \regfield{EFUSE\_DIS\_DOWNLOAD\_MODE}{1}{0}{{0}}%
    \reglabel{Reset}\regnewline%
    \begin{regdesc}\begin{reglist}
    \label{fielddesc:EFUSEDISDOWNLOADMODE}\item [EFUSE\_DIS\_DOWNLOAD\_MODE] 表示是否关闭所有 Download 模式。\\
    1：关闭\\
    0：不关闭\\
    (RO)
    \label{fielddesc:EFUSEDISDIRECTBOOT}\item [EFUSE\_DIS\_DIRECT\_BOOT] 表示是否使能 direct boot 模式。\\
    1：不使能\\
    0：使能\\
    (RO)
    \label{fielddesc:EFUSEDISUSBSERIALJTAGROMPRINT}\item [EFUSE\_DIS\_USB\_SERIAL\_JTAG\_ROM\_PRINT] 表示是否使能 USB-Serial-JTAG 打印。\\
    1：不使能\\
    0：使能\\
    (RO)
    \label{fielddesc:EFUSEDISUSBSERIALJTAGDOWNLOADMODE}\item [EFUSE\_DIS\_USB\_SERIAL\_JTAG\_DOWNLOAD\_MODE] 表示是否使能 USB-Serial-JTAG 下载功能。\\
    1：不使能\\
    0：使能\\
    (RO)
    \label{fielddesc:EFUSEENABLESECURITYDOWNLOAD}\item [EFUSE\_ENABLE\_SECURITY\_DOWNLOAD] 表示是否使能安全下载。\\
    1：使能\\
    0：不使能\\
    (RO)
    \label{fielddesc:EFUSEUARTPRINTCONTROL}\item [EFUSE\_UART\_PRINT\_CONTROL] 表示 UART 打印类型。\\
    0：强制打印\\
    1：GPIO8 低电平复位时使能打印\\
    2：GPIO8 高电平复位时使能打印\\
    3：强制关闭打印\\
    (RO)
    \label{fielddesc:EFUSEFORCESENDRESUME}\item [EFUSE\_FORCE\_SEND\_RESUME] 表示是否强制 ROM 代码在 SPI 启动过程中发送恢复指令。 \\
    1：强制\\
    0：不强制\\
    (RO)
    \label{fielddesc:EFUSESECUREVERSION}\item [EFUSE\_SECURE\_VERSION] 表示安全版本，用于 ESP-IDF 防回滚功能。(RO)

    \item [见下页……]
    \end{reglist}\end{regdesc}
    \end{register}

    \addtocounter{Regfloat}{-1}
    \begin{register}{H}{EFUSE\_RD\_REPEAT\_DATA2\_REG}{0x0038}
    \begin{regdesc}\begin{reglist}
    \item [接上页……]

    \label{fielddesc:EFUSESECUREBOOTDISABLEFASTWAKE}\item [EFUSE\_SECURE\_BOOT\_DISABLE\_FAST\_WAKE] 表示安全启动使能时是否使能 FAST\_VERIFY\_ON\_WAKE。\\
    1：不使能\\
    0：使能\\
    (RO)
    \label{fielddesc:EFUSEHYSENPAD}\item [EFUSE\_HYS\_EN\_PAD] 表示是否启用对应 PAD 的迟滞功能。\\
    1：使能\\
    0：不使能\\
    (RO)
    \label{fielddesc:EFUSEXTSDPACLKENABLE}\item [EFUSE\_XTS\_DPA\_CLK\_ENABLE] 表示是否启用抗 DPA 攻击的时钟功能。\\
    1：使能\\
    0：不使能\\
    (RO)
    \label{fielddesc:EFUSEXTSDPAPSEUDOLEVEL}\item [EFUSE\_XTS\_DPA\_PSEUDO\_LEVEL] 表示伪轮机制的抗旁路攻击能力等级。\\
    3：高等级\\
    2：中等级\\
    1：低等级\\
    0：由寄存器配置决定\\
    (RO)
    \label{fielddesc:EFUSEDISWIFI6}\item [EFUSE\_DIS\_WIFI6] 表示是否使能 Wi-Fi 6。\\
    1：不使能\\
    0：使能\\
    (RO)
    \label{fielddesc:EFUSEECDSADISABLEP192}\item [EFUSE\_ECDSA\_DISABLE\_P192] 表示在 ECDSA\_DS 中是否使能 P192 曲线。\\
    1：不使能\\
    0：使能 \\
    (RO)
    \label{fielddesc:EFUSEECCFORCECONSTTIME}\item [EFUSE\_ECC\_FORCE\_CONST\_TIME] 表示是否强制 ECC 使用固定时间的计算模式。\\
    1：使能 \\
    0：不使能 \\
    (RO)
    \end{reglist}\end{regdesc}
    \end{register}


    \begin{register}{H}{EFUSE\_RD\_REPEAT\_DATA3\_REG}{0x003C}\label{regdesc:EFUSERDREPEATDATA3REG}
    \regfield{EFUSE\_RD\_REPEAT\_DATA3}{32}{0}{{0x000000}}%
    \reglabel{Reset}\regnewline%
    \begin{regdesc}\begin{reglist}
    \label{fielddesc:EFUSERDREPEATDATA3}\item [EFUSE\_RD\_REPEAT\_DATA3] 保留 (RO)
    \end{reglist}\end{regdesc}
    \end{register}


    \begin{register}{H}{EFUSE\_RD\_REPEAT\_DATA4\_REG}{0x0040}\label{regdesc:EFUSERDREPEATDATA4REG}
    \regfield{EFUSE\_RD\_REPEAT\_DATA4}{32}{0}{{0x000000}}%
    \reglabel{Reset}\regnewline%
    \begin{regdesc}\begin{reglist}
    \label{fielddesc:EFUSERDREPEATDATA4}\item [EFUSE\_RD\_REPEAT\_DATA4] 保留 (RO)
    \end{reglist}\end{regdesc}
    \end{register}
    

    \begin{register}{H}{EFUSE\_RD\_MAC\_SYS0\_REG}{0x0044}\label{regdesc:EFUSERDMACSYS0REG}
    \regfield{EFUSE\_MAC\_0}{32}{0}{{0x000000}}%
    \reglabel{Reset}\regnewline%
    \begin{regdesc}\begin{reglist}
    \label{fielddesc:EFUSEMAC0}\item [EFUSE\_MAC\_0] 表示 MAC 地址低 32 位。(RO)
    \end{reglist}\end{regdesc}
    \end{register}


    \begin{register}{H}{EFUSE\_RD\_MAC\_SYS1\_REG}{0x0048}\label{regdesc:EFUSERDMACSYS1REG}
    \regfield{EFUSE\_MAC\_EXT}{16}{16}{{0x00}}%
    \regfield{EFUSE\_MAC\_1}{16}{0}{{0x00}}%
    \reglabel{Reset}\regnewline%
    \begin{regdesc}\begin{reglist}
    \label{fielddesc:EFUSEMAC1}\item [EFUSE\_MAC\_1] 表示 MAC 地址高 16 位。 (RO)
    \label{fielddesc:EFUSEMACEXT}\item [EFUSE\_MAC\_EXT] 表示 MAC 地址扩展位。 (RO)
    \end{reglist}\end{regdesc}
    \end{register}


    \begin{register}{H}{EFUSE\_RD\_MAC\_SYS2\_REG}{0x004C}\label{regdesc:EFUSERDMACSYS2REG}
    \regfield{EFUSE\_MAC\_RESERVED\_1}{18}{14}{{0x000}}%
    \regfield{EFUSE\_MAC\_RESERVED\_0}{14}{0}{{0x00}}%
    \reglabel{Reset}\regnewline%
    \begin{regdesc}\begin{reglist}
    \label{fielddesc:EFUSEMACRESERVED0}\item [EFUSE\_MAC\_RESERVED\_0] 保留 (RO)
    \label{fielddesc:EFUSEMACRESERVED1}\item [EFUSE\_MAC\_RESERVED\_1] 保留 (RO)
    \end{reglist}\end{regdesc}
    \end{register}


    \begin{register}{H}{EFUSE\_RD\_MAC\_SYS3\_REG}{0x0050}\label{regdesc:EFUSERDMACSYS3REG}
    \regfield{EFUSE\_SYS\_DATA\_PART0\_0}{14}{18}{{0x00}}%
    \regfield{EFUSE\_MAC\_RESERVED\_2}{18}{0}{{0x000}}%
    \reglabel{Reset}\regnewline%
    \begin{regdesc}\begin{reglist}
    \label{fielddesc:EFUSEMACRESERVED2}\item [EFUSE\_MAC\_RESERVED\_2] 保留 (RO)
    \label{fielddesc:EFUSESYSDATAPART00}\item [EFUSE\_SYS\_DATA\_PART0\_0] 表示系统数据第 0 部分的第 1 个 14 位内容。(RO)
    \end{reglist}\end{regdesc}
    \end{register}


    \begin{register}{H}{EFUSE\_RD\_MAC\_SYS4\_REG}{0x0054}\label{regdesc:EFUSERDMACSYS4REG}
    \regfield{EFUSE\_SYS\_DATA\_PART0\_1}{32}{0}{{0x000000}}%
    \reglabel{Reset}\regnewline%
    \begin{regdesc}\begin{reglist}
    \label{fielddesc:EFUSESYSDATAPART01}\item [EFUSE\_SYS\_DATA\_PART0\_1] 表示系统数据第 0 部分的第 2 个 32 位内容。(RO)
    \end{reglist}\end{regdesc}
    \end{register}


    \begin{register}{H}{EFUSE\_RD\_MAC\_SYS5\_REG}{0x0058}\label{regdesc:EFUSERDMACSYS5REG}
    \regfield{EFUSE\_SYS\_DATA\_PART0\_2}{32}{0}{{0x000000}}%
    \reglabel{Reset}\regnewline%
    \begin{regdesc}\begin{reglist}
    \label{fielddesc:EFUSESYSDATAPART02}\item [EFUSE\_SYS\_DATA\_PART0\_2] 表示系统数据第 0 部分的第 3 个 32 位内容。 (RO)
    \end{reglist}\end{regdesc}
    \end{register}


    \begin{register}{H}{EFUSE\_RD\_SYS\_PART1\_DATA\regindex{n}\_REG (\regindex{n}: 0-7)}{0x005C+0x4*\regindex{n}}\label{regdesc:EFUSERDSYSPART1DATANREG}
    \regfield{EFUSE\_SYS\_DATA\_PART1\_0}{32}{0}{{0x000000}}%
    \reglabel{Reset}\regnewline%
    \begin{regdesc}\begin{reglist}
    \label{fielddesc:EFUSESYSDATAPART1N}\item [EFUSE\_SYS\_DATA\_PART1\_\regindex{n}] 表示系统数据第 1 部分的第 \regindex{n} 个 32 位内容。(RO)
    \end{reglist}\end{regdesc}
    \end{register}


    \begin{register}{H}{EFUSE\_RD\_USR\_DATA\regindex{n}\_REG (\regindex{n}: 0-7)}{0x007C+0x4*\regindex{n}}\label{regdesc:EFUSERDUSRDATANREG}
    \regfield{EFUSE\_USR\_DATA0}{32}{0}{{0x000000}}%
    \reglabel{Reset}\regnewline%
    \begin{regdesc}\begin{reglist}
    \label{fielddesc:EFUSEUSRDATAN}\item [EFUSE\_USR\_DATA\regindex{n}] 表示 BLOCK3 (user) 第 \regindex{n} 个 32 位内容。(RO)
    \end{reglist}\end{regdesc}
    \end{register}


    \begin{register}{H}{EFUSE\_RD\_KEY0\_DATA\regindex{n}\_REG (\regindex{n}: 0-7)}{0x009C+0x4*\regindex{n}}\label{regdesc:EFUSERDKEY0DATANREG}
    \regfield{EFUSE\_KEY0\_DATA0}{32}{0}{{0x000000}}%
    \reglabel{Reset}\regnewline%
    \begin{regdesc}\begin{reglist}
    \label{fielddesc:EFUSEKEY0DATAN}\item [EFUSE\_KEY0\_DATA\regindex{n}] 表示 KEY0 第 \regindex{n} 个 32 位内容。(RO)
    \end{reglist}\end{regdesc}
    \end{register}


    \begin{register}{H}{EFUSE\_RD\_KEY1\_DATA\regindex{n}\_REG (\regindex{n}: 0-7)}{0x00BC+0x4*\regindex{n}}\label{regdesc:EFUSERDKEY1DATANREG}
    \regfield{EFUSE\_KEY1\_DATA0}{32}{0}{{0x000000}}%
    \reglabel{Reset}\regnewline%
    \begin{regdesc}\begin{reglist}
    \label{fielddesc:EFUSEKEY1DATAN}\item [EFUSE\_KEY1\_DATA\regindex{n}] 表示 KEY1 第 \regindex{n} 个 32 位内容。 (RO)
    \end{reglist}\end{regdesc}
    \end{register}


    \begin{register}{H}{EFUSE\_RD\_KEY2\_DATA\regindex{n}\_REG (\regindex{n}: 0-7)}{0x00DC+0x4*\regindex{n}}\label{regdesc:EFUSERDKEY2DATANREG}
    \regfield{EFUSE\_KEY2\_DATA0}{32}{0}{{0x000000}}%
    \reglabel{Reset}\regnewline%
    \begin{regdesc}\begin{reglist}
    \label{fielddesc:EFUSEKEY2DATAN}\item [EFUSE\_KEY2\_DATA\regindex{n}] 表示 KEY2 第 \regindex{n} 个 32 位内容。 (RO)
    \end{reglist}\end{regdesc}
    \end{register}


    \begin{register}{H}{EFUSE\_RD\_KEY3\_DATA\regindex{n}\_REG (\regindex{n}: 0-7)}{0x00FC+0x4*\regindex{n}}\label{regdesc:EFUSERDKEY3DATANREG}
    \regfield{EFUSE\_KEY3\_DATA0}{32}{0}{{0x000000}}%
    \reglabel{Reset}\regnewline%
    \begin{regdesc}\begin{reglist}
    \label{fielddesc:EFUSEKEY3DATAN}\item [EFUSE\_KEY3\_DATA\regindex{n}] 表示 KEY3 第 \regindex{n} 个 32 位内容。(RO)
    \end{reglist}\end{regdesc}
    \end{register}


    \begin{register}{H}{EFUSE\_RD\_KEY4\_DATA\regindex{n}\_REG (\regindex{n}: 0-7)}{0x011C+0x4*\regindex{n}}\label{regdesc:EFUSERDKEY4DATANREG}
    \regfield{EFUSE\_KEY4\_DATA0}{32}{0}{{0x000000}}%
    \reglabel{Reset}\regnewline%
    \begin{regdesc}\begin{reglist}
    \label{fielddesc:EFUSEKEY4DATAN}\item [EFUSE\_KEY4\_DATA\regindex{n}] 表示 KEY4 第 \regindex{n} 个 32 位内容。(RO)
    \end{reglist}\end{regdesc}
    \end{register}


    \begin{register}{H}{EFUSE\_RD\_KEY5\_DATA\regindex{n}\_REG (\regindex{n}: 0-7)}{0x013C+0x4*\regindex{n}}\label{regdesc:EFUSERDKEY5DATANREG}
    \regfield{EFUSE\_KEY5\_DATA0}{32}{0}{{0x000000}}%
    \reglabel{Reset}\regnewline%
    \begin{regdesc}\begin{reglist}
    \label{fielddesc:EFUSEKEY5DATAN}\item [EFUSE\_KEY5\_DATA\regindex{n}] 表示 KEY5 第 \regindex{n} 个 32 位内容。(RO)
    \end{reglist}\end{regdesc}
    \end{register}


    \begin{register}{H}{EFUSE\_RD\_SYS\_PART2\_DATA\regindex{n}\_REG (\regindex{n}: 0-7)}{0x015C+0x4*\regindex{n}}\label{regdesc:EFUSERDSYSPART2DATANREG}
    \regfield{EFUSE\_SYS\_DATA\_PART2\_0}{32}{0}{{0x000000}}%
    \reglabel{Reset}\regnewline%
    \begin{regdesc}\begin{reglist}
    \label{fielddesc:EFUSESYSDATAPART2N}\item [EFUSE\_SYS\_DATA\_PART2\_\regindex{n}] 表示系统数据第 2 部分的第 \regindex{n} 个 32 位内容。(RO)
    \end{reglist}\end{regdesc}
    \end{register}


    \begin{register}{H}{EFUSE\_RD\_REPEAT\_DATA\_ERR0\_REG}{0x017C}\label{regdesc:EFUSERDREPEATDATAERR0REG}
    \regfield{(reserved)}{3}{29}{000}%
    \regfield{EFUSE\_SECURE\_BOOT\_KEY\_REVOKE\_2\_ERR}{1}{28}{{0}}%
    \regfield{EFUSE\_SECURE\_BOOT\_KEY\_REVOKE\_1\_ERR}{1}{27}{{0}}%
    \regfield{EFUSE\_SECURE\_BOOT\_KEY\_REVOKE\_0\_ERR}{1}{26}{{0}}%
    \regfield{EFUSE\_SPI\_BOOT\_CRYPT\_CNT\_ERR}{3}{23}{{0x0}}%
    \regfield{EFUSE\_WDT\_DELAY\_SEL\_ERR}{2}{21}{{0x0}}%
    \regfield{EFUSE\_VDD\_SPI\_AS\_GPIO\_ERR}{1}{20}{{0}}%
    \regfield{EFUSE\_USB\_EXCHG\_PINS\_ERR}{1}{19}{{0}}%
    \regfield{EFUSE\_USB\_DREFL\_ERR}{2}{17}{{0x0}}%
    \regfield{EFUSE\_USB\_DREFH\_ERR}{2}{15}{{0x0}}%
    \regfield{EFUSE\_DIS\_DOWNLOAD\_MANUAL\_ENCRYPT\_ERR}{1}{14}{{0}}%
    \regfield{EFUSE\_DIS\_PAD\_JTAG\_ERR}{1}{13}{{0}}%
    \regfield{EFUSE\_JTAG\_SEL\_ENABLE\_ERR}{1}{12}{{0}}%
    \regfield{EFUSE\_SPI\_DOWNLOAD\_MSPI\_DIS\_ERR}{1}{11}{{0}}%
    \regfield{EFUSE\_DIS\_FORCE\_DOWNLOAD\_ERR}{1}{10}{{0}}%
    \regfield{EFUSE\_DIS\_USB\_SERIAL\_JTAG\_ERR}{1}{9}{{0}}%
    \regfield{EFUSE\_DIS\_USB\_JTAG\_ERR}{1}{8}{{0}}%
    \regfield{EFUSE\_DIS\_ICACHE\_ERR}{1}{7}{{0}}%
    \regfield{EFUSE\_RD\_DIS\_ERR}{7}{0}{{0x0}}%
    \reglabel{Reset}\regnewline%
    \begin{regdesc}\begin{reglist}
    \label{fielddesc:EFUSERDDISERR}\item [EFUSE\_RD\_DIS\_ERR] 若该位为 1，则表示 \hyperref[fielddesc:EFUSERDDIS]{EFUSE\_RD\_DIS} 烧写错误。(RO)
    \label{fielddesc:EFUSEDISICACHEERR}\item [EFUSE\_DIS\_ICACHE\_ERR] 若该位为 1，则表示 \hyperref[fielddesc:EFUSEDISICACHE]{EFUSE\_DIS\_ICACHE} 烧写错误。(RO)
    \label{fielddesc:EFUSEDISUSBJTAGERR}\item [EFUSE\_DIS\_USB\_JTAG\_ERR] 若该位为 1，则表示 \hyperref[fielddesc:EFUSEDISUSBJTAG]{EFUSE\_DIS\_USB\_JTAG} 烧写错误。(RO)
    \label{fielddesc:EFUSEDISUSBSERIALJTAGERR}\item [EFUSE\_DIS\_USB\_SERIAL\_JTAG\_ERR] 若该位为 1，则表示 \hyperref[fielddesc:EFUSEDISUSBSERIALJTAG]{EFUSE\_DIS\_USB\_SERIAL\_JTAG} 烧写错误。(RO)
    \label{fielddesc:EFUSEDISFORCEDOWNLOADERR}\item [EFUSE\_DIS\_FORCE\_DOWNLOAD\_ERR] 若该位为 1，则表示 \hyperref[fielddesc:EFUSEDISFORCEDOWNLOAD]{EFUSE\_DIS\_FORCE\_DOWNLOAD} 烧写错误。(RO)
    \label{fielddesc:EFUSESPIDOWNLOADMSPIDISERR}\item [EFUSE\_SPI\_DOWNLOAD\_MSPI\_DIS\_ERR] 若该位为 1，则表示 \hyperref[fielddesc:EFUSESPIDOWNLOADMSPIDIS]{EFUSE\_SPI\_DOWNLOAD\_MSPI\_DIS} 烧写错误。(RO)
    \label{fielddesc:EFUSEJTAGSELENABLEERR}\item [EFUSE\_JTAG\_SEL\_ENABLE\_ERR] 若该位为 1，则表示 \hyperref[fielddesc:EFUSEJTAGSELENABLE]{EFUSE\_JTAG\_SEL\_ENABLE} 烧写错误。(RO)
    \label{fielddesc:EFUSEDISPADJTAGERR}\item [EFUSE\_DIS\_PAD\_JTAG\_ERR] 若该位为 1，则表示 \hyperref[fielddesc:EFUSEDISPADJTAG]{EFUSE\_DIS\_PAD\_JTAG} 烧写错误。(RO)
    \label{fielddesc:EFUSEDISDOWNLOADMANUALENCRYPTERR}\item [EFUSE\_DIS\_DOWNLOAD\_MANUAL\_ENCRYPT\_ERR] 若该位为 1，则表示 \hyperref[fielddesc:EFUSEDISDOWNLOADMANUALENCRYPT]{EFUSE\_DIS\_DOWNLOAD\_MANUAL\_ENCRYPT} 烧写错误。(RO)
    \label{fielddesc:EFUSEUSBDREFHERR}\item [EFUSE\_USB\_DREFH\_ERR] 若该位为 1，则表示 \hyperref[fielddesc:EFUSEUSBDREFH]{EFUSE\_USB\_DREFH} 烧写错误。(RO)
    \label{fielddesc:EFUSEUSBDREFLERR}\item [EFUSE\_USB\_DREFL\_ERR] 若该位为 1，则表示 \hyperref[fielddesc:EFUSEUSBDREFL]{EFUSE\_USB\_DREFL} 烧写错误。(RO)
    \label{fielddesc:EFUSEUSBEXCHGPINSERR}\item [EFUSE\_USB\_EXCHG\_PINS\_ERR] 若该位为 1，则表示 \hyperref[fielddesc:EFUSEUSBEXCHGPINS]{EFUSE\_USB\_EXCHG\_PINS} 烧写错误。(RO)
    \label{fielddesc:EFUSEVDDSPIASGPIOERR}\item [EFUSE\_VDD\_SPI\_AS\_GPIO\_ERR] 若该位为 1，则表示 \hyperref[fielddesc:EFUSEVDDSPIASGPIO]{EFUSE\_VDD\_SPI\_AS\_GPIO} 烧写错误。(RO)
    \label{fielddesc:EFUSEWDTDELAYSELERR}\item [EFUSE\_WDT\_DELAY\_SEL\_ERR] 若该位为 1，则表示 \hyperref[fielddesc:EFUSEWDTDELAYSEL]{EFUSE\_WDT\_DELAY\_SEL} 烧写错误。(RO)
    \label{fielddesc:EFUSESPIBOOTCRYPTCNTERR}\item [EFUSE\_SPI\_BOOT\_CRYPT\_CNT\_ERR] 若该位为 1，则表示 \hyperref[fielddesc:EFUSESPIBOOTCRYPTCNT]{EFUSE\_SPI\_BOOT\_CRYPT\_CNT} 烧写错误。(RO)

    \item [见下页……]
    \end{reglist}\end{regdesc}
    \end{register}

    \addtocounter{Regfloat}{-1}
    \begin{register}{H}{EFUSE\_RD\_REPEAT\_DATA\_ERR0\_REG}{0x017C}
    \begin{regdesc}\begin{reglist}
    \item [接上页……]

    \label{fielddesc:EFUSESECUREBOOTKEYREVOKE0ERR}\item [EFUSE\_SECURE\_BOOT\_KEY\_REVOKE\_0\_ERR] 若该位为 1，则表示 \hyperref[fielddesc:EFUSESECUREBOOTKEYREVOKE0]{EFUSE\_SECURE\_BOOT\_KEY\_REVOKE\_0} 烧写错误。(RO)
    \label{fielddesc:EFUSESECUREBOOTKEYREVOKE1ERR}\item [EFUSE\_SECURE\_BOOT\_KEY\_REVOKE\_1\_ERR] 若该位为 1，则表示 \hyperref[fielddesc:EFUSESECUREBOOTKEYREVOKE1]{EFUSE\_SECURE\_BOOT\_KEY\_REVOKE\_1} 烧写错误。(RO)
    \label{fielddesc:EFUSESECUREBOOTKEYREVOKE2ERR}\item [EFUSE\_SECURE\_BOOT\_KEY\_REVOKE\_2\_ERR] 若该位为 1，则表示 \hyperref[fielddesc:EFUSESECUREBOOTKEYREVOKE2]{EFUSE\_SECURE\_BOOT\_KEY\_REVOKE\_2} 烧写错误。(RO)
    \end{reglist}\end{regdesc}
    \end{register}


    \begin{register}{H}{EFUSE\_RD\_REPEAT\_DATA\_ERR1\_REG}{0x0180}\label{regdesc:EFUSERDREPEATDATAERR1REG}
    \regfield{EFUSE\_FLASH\_TPUW\_ERR}{4}{28}{{0x0}}%
    \regfield{EFUSE\_SECURE\_BOOT\_AGGRESSIVE\_REVOKE\_ERR}{1}{27}{{0}}%
    \regfield{EFUSE\_SECURE\_BOOT\_EN\_ERR}{1}{26}{{0}}%
    \regfield{EFUSE\_SEC\_DPA\_LEVEL\_ERR}{2}{24}{{0x0}}%
    \regfield{EFUSE\_KEY\_PURPOSE\_5\_ERR}{4}{20}{{0x0}}%
    \regfield{EFUSE\_KEY\_PURPOSE\_4\_ERR}{4}{16}{{0x0}}%
    \regfield{EFUSE\_KEY\_PURPOSE\_3\_ERR}{4}{12}{{0x0}}%
    \regfield{EFUSE\_KEY\_PURPOSE\_2\_ERR}{4}{8}{{0x0}}%
    \regfield{EFUSE\_KEY\_PURPOSE\_1\_ERR}{4}{4}{{0x0}}%
    \regfield{EFUSE\_KEY\_PURPOSE\_0\_ERR}{4}{0}{{0x0}}%
    \reglabel{Reset}\regnewline%
    \begin{regdesc}\begin{reglist}
    \label{fielddesc:EFUSEKEYPURPOSE0ERR}\item [EFUSE\_KEY\_PURPOSE\_0\_ERR] 若该位为 1，则表示 \hyperref[fielddesc:EFUSEKEYPURPOSE0]{EFUSE\_KEY\_PURPOSE\_0} 烧写错误。(RO)
    \label{fielddesc:EFUSEKEYPURPOSE1ERR}\item [EFUSE\_KEY\_PURPOSE\_1\_ERR] 若该位为 1，则表示 \hyperref[fielddesc:EFUSEKEYPURPOSE1]{EFUSE\_KEY\_PURPOSE\_1} 烧写错误。(RO)
    \label{fielddesc:EFUSEKEYPURPOSE2ERR}\item [EFUSE\_KEY\_PURPOSE\_2\_ERR] 若该位为 1，则表示 \hyperref[fielddesc:EFUSEKEYPURPOSE2]{EFUSE\_KEY\_PURPOSE\_2} 烧写错误。(RO)
    \label{fielddesc:EFUSEKEYPURPOSE3ERR}\item [EFUSE\_KEY\_PURPOSE\_3\_ERR] 若该位为 1，则表示 \hyperref[fielddesc:EFUSEKEYPURPOSE3]{EFUSE\_KEY\_PURPOSE\_3} 烧写错误。(RO)
    \label{fielddesc:EFUSEKEYPURPOSE4ERR}\item [EFUSE\_KEY\_PURPOSE\_4\_ERR] 若该位为 1，则表示 \hyperref[fielddesc:EFUSEKEYPURPOSE4]{EFUSE\_KEY\_PURPOSE\_4} 烧写错误。(RO)
    \label{fielddesc:EFUSEKEYPURPOSE5ERR}\item [EFUSE\_KEY\_PURPOSE\_5\_ERR] 若该位为 1，则表示 \hyperref[fielddesc:EFUSEKEYPURPOSE5]{EFUSE\_KEY\_PURPOSE\_5} 烧写错误。(RO)
    \label{fielddesc:EFUSESECDPALEVELERR}\item [EFUSE\_SEC\_DPA\_LEVEL\_ERR] 若该位为 1，则表示 \hyperref[fielddesc:EFUSESECDPALEVEL]{EFUSE\_SEC\_DPA\_LEVEL} 烧写错误。(RO)
    \label{fielddesc:EFUSESECUREBOOTENERR}\item [EFUSE\_SECURE\_BOOT\_EN\_ERR] 若该位为 1，则表示 \hyperref[fielddesc:EFUSESECUREBOOTEN]{EFUSE\_SECURE\_BOOT\_EN} 烧写错误。(RO)
    \label{fielddesc:EFUSESECUREBOOTAGGRESSIVEREVOKEERR}\item [EFUSE\_SECURE\_BOOT\_AGGRESSIVE\_REVOKE\_ERR] 若该位为 1，则表示 \hyperref[fielddesc:EFUSESECUREBOOTAGGRESSIVEREVOKE]{EFUSE\_SECURE\_BOOT\_AGGRESSIVE\_REVOKE} 烧写错误。(RO)
    \label{fielddesc:EFUSEFLASHTPUWERR}\item [EFUSE\_FLASH\_TPUW\_ERR] 若该位为 1，则表示 \hyperref[fielddesc:EFUSEFLASHTPUW]{EFUSE\_FLASH\_TPUW} 烧写错误。(RO)
    \end{reglist}\end{regdesc}
    \end{register}


    \begin{register}{H}{EFUSE\_RD\_REPEAT\_DATA\_ERR2\_REG}{0x0184}\label{regdesc:EFUSERDREPEATDATAERR2REG}
    \regfield{EFUSE\_ECC\_FORCE\_CONST\_TIME\_ERR}{1}{31}{{0}}%
    \regfield{EFUSE\_ECDSA\_DISABLE\_P192\_ERR}{1}{30}{{0}}%
    \regfield{EFUSE\_DIS\_WIFI6\_ERR}{1}{29}{{0}}%
    \regfield{EFUSE\_XTS\_DPA\_PSEUDO\_LEVEL\_ERR}{2}{27}{{0x0}}%
    \regfield{EFUSE\_XTS\_DPA\_CLK\_ENABLE\_ERR}{1}{26}{{0}}%
    \regfield{EFUSE\_HYS\_EN\_PAD\_ERR}{1}{25}{{0}}%
    \regfield{EFUSE\_SECURE\_BOOT\_DISABLE\_FAST\_WAKE\_ERR}{1}{24}{{0}}%
    \regfield{EFUSE\_SECURE\_VERSION\_ERR}{16}{8}{{0x00}}%
    \regfield{\rotatebox{10}{EFUSE\_FORCE\_SEND\_RESUME\_ERR}}{1}{7}{{0}}%
    \regfield{\rotatebox{10}{EFUSE\_UART\_PRINT\_CONTROL\_ERR}}{2}{5}{{0x0}}%
    \regfield{\rotatebox{10}{EFUSE\_ENABLE\_SECURITY\_DOWNLOAD\_ERR}}{1}{4}{{0}}%
    \regfield{\rotatebox{10}{EFUSE\_DIS\_USB\_SERIAL\_JTAG\_DOWNLOAD\_MODE\_ERR}}{1}{3}{{0}}%
    \regfield{\rotatebox{10}{EFUSE\_DIS\_USB\_SERIAL\_JTAG\_ROM\_PRINT\_ERR}}{1}{2}{{0}}%
    \regfield{\rotatebox{10}{EFUSE\_DIS\_DIRECT\_BOOT\_ERR}}{1}{1}{{0}}%
    \regfield{\rotatebox{10}{EFUSE\_DIS\_DOWNLOAD\_MODE\_ERR}}{1}{0}{{0}}%
    \reglabel{Reset}\regnewline%
    \begin{regdesc}\begin{reglist}
    \label{fielddesc:EFUSEDISDOWNLOADMODEERR}\item [EFUSE\_DIS\_DOWNLOAD\_MODE\_ERR] 若该位为 1，则表示 \hyperref[fielddesc:EFUSEDISDOWNLOADMODE]{EFUSE\_DIS\_DOWNLOAD\_MODE} 烧写错误。(RO)
    \label{fielddesc:EFUSEDISDIRECTBOOTERR}\item [EFUSE\_DIS\_DIRECT\_BOOT\_ERR] 若该位为 1，则表示 \hyperref[fielddesc:EFUSEDISDIRECTBOOT]{EFUSE\_DIS\_DIRECT\_BOOT} 烧写错误。(RO)
    \label{fielddesc:EFUSEDISUSBSERIALJTAGROMPRINTERR}\item [EFUSE\_DIS\_USB\_SERIAL\_JTAG\_ROM\_PRINT\_ERR] 若该位为 1，则表示 \hyperref[fielddesc:EFUSEDISUSBSERIALJTAGROMPRINT]{EFUSE\_DIS\_USB\_SERIAL\_JTAG\_ROM\_PRINT} 烧写错误。(RO)
    \label{fielddesc:EFUSEDISUSBSERIALJTAGDOWNLOADMODEERR}\item [EFUSE\_DIS\_USB\_SERIAL\_JTAG\_DOWNLOAD\_MODE\_ERR] 若该位为 1，则表示 \hyperref[fielddesc:EFUSEDISUSBSERIALJTAGDOWNLOADMODE]{EFUSE\_DIS\_USB\_SERIAL\_JTAG\_DOWNLOAD\_MODE} 烧写错误。(RO)
    \label{fielddesc:EFUSEENABLESECURITYDOWNLOADERR}\item [EFUSE\_ENABLE\_SECURITY\_DOWNLOAD\_ERR] 若该位为 1，则表示 \hyperref[fielddesc:EFUSEENABLESECURITYDOWNLOAD]{EFUSE\_ENABLE\_SECURITY\_DOWNLOAD} 烧写错误。(RO)
    \label{fielddesc:EFUSEUARTPRINTCONTROLERR}\item [EFUSE\_UART\_PRINT\_CONTROL\_ERR] 若该位为 1，则表示 \hyperref[fielddesc:EFUSEUARTPRINTCONTROL]{EFUSE\_UART\_PRINT\_CONTROL} 烧写错误。(RO)
    \label{fielddesc:EFUSEFORCESENDRESUMEERR}\item [EFUSE\_FORCE\_SEND\_RESUME\_ERR] 若该位为 1，则表示 \hyperref[fielddesc:EFUSEFORCESENDRESUME]{EFUSE\_FORCE\_SEND\_RESUME} 烧写错误。(RO)
    \label{fielddesc:EFUSESECUREVERSIONERR}\item [EFUSE\_SECURE\_VERSION\_ERR] 若该位为 1，则表示 \hyperref[fielddesc:EFUSESECUREVERSION]{EFUSE\_SECURE\_VERSION} 烧写错误。(RO)
    \label{fielddesc:EFUSESECUREBOOTDISABLEFASTWAKEERR}\item [EFUSE\_SECURE\_BOOT\_DISABLE\_FAST\_WAKE\_ERR] 若该位为 1，则表示 \hyperref[fielddesc:EFUSESECUREBOOTDISABLEFASTWAKE]{EFUSE\_SECURE\_BOOT\_DISABLE\_FAST\_WAKE} 烧写错误。(RO)
    \label{fielddesc:EFUSEHYSENPADERR}\item [EFUSE\_HYS\_EN\_PAD\_ERR] 若该位为 1，则表示 \hyperref[fielddesc:EFUSEHYSENPAD]{EFUSE\_HYS\_EN\_PAD} 烧写错误。(RO)
    \label{fielddesc:EFUSEXTSDPACLKENABLEERR}\item [EFUSE\_XTS\_DPA\_CLK\_ENABLE\_ERR] 若该位为 1，则表示 \hyperref[fielddesc:EFUSEXTSDPACLKENABLE]{EFUSE\_XTS\_DPA\_CLK\_ENABLE} 烧写错误。(RO)
    \label{fielddesc:EFUSEXTSDPAPSEUDOLEVELERR}\item [EFUSE\_XTS\_DPA\_PSEUDO\_LEVEL\_ERR] 若该位为 1，则表示 \hyperref[fielddesc:EFUSEXTSDPAPSEUDOLEVEL]{EFUSE\_XTS\_DPA\_PSEUDO\_LEVEL}. (RO)
    \label{fielddesc:EFUSEDISWIFI6ERR}\item [EFUSE\_DIS\_WIFI6\_ERR] 若该位为 1，则表示 \hyperref[fielddesc:EFUSEDISWIFI6]{EFUSE\_DIS\_WIFI6} 烧写错误。(RO)

    \item [见下页……]
    \end{reglist}\end{regdesc}
    \end{register}

    \addtocounter{Regfloat}{-1}
    \begin{register}{H}{EFUSE\_RD\_REPEAT\_DATA\_ERR2\_REG}{0x0184}
    \begin{regdesc}\begin{reglist}
    \item [接上页……]

    \label{fielddesc:EFUSEECDSADISABLEP192ERR}\item [EFUSE\_ECDSA\_DISABLE\_P192\_ERR] 若该位为 1，则表示 \hyperref[fielddesc:EFUSEECDSADISABLEP192]{EFUSE\_ECDSA\_DISABLE\_P192} 烧写错误。(RO)
    \label{fielddesc:EFUSEECCFORCECONSTTIMEERR}\item [EFUSE\_ECC\_FORCE\_CONST\_TIME\_ERR] 若该位为 1，则表示 \hyperref[fielddesc:EFUSEECCFORCECONSTTIME]{EFUSE\_ECC\_FORCE\_CONST\_TIME} 烧写错误。(RO)
    \end{reglist}\end{regdesc}
    \end{register}
    

    \begin{register}{H}{EFUSE\_RD\_REPEAT\_DATA\_ERR3\_REG}{0x0188}\label{regdesc:EFUSERDREPEATDATAERR3REG}
    \regfield{EFUSE\_RD\_REPEAT\_DATA\_ERR3}{32}{0}{{0x000000}}%
    \reglabel{Reset}\regnewline%
    \begin{regdesc}\begin{reglist}
    \label{fielddesc:EFUSERDREPEATDATAERR3}\item [EFUSE\_RD\_REPEAT\_DATA\_ERR3] 保留 (RO)
    \end{reglist}\end{regdesc}
    \end{register}


    \begin{register}{H}{EFUSE\_RD\_REPEAT\_DATA\_ERR4\_REG}{0x018C}\label{regdesc:EFUSERDREPEATDATAERR4REG}
    \regfield{EFUSE\_RD\_REPEAT\_DATA\_ERR4}{32}{0}{{0x000000}}%
    \reglabel{Reset}\regnewline%
    \begin{regdesc}\begin{reglist}
    \label{fielddesc:EFUSERDREPEATDATAERR4}\item [EFUSE\_RD\_REPEAT\_DATA\_ERR4] 保留 (RO)
    \end{reglist}\end{regdesc}
    \end{register}


    \begin{register}{H}{EFUSE\_RD\_RS\_DATA\_ERR0\_REG}{0x0190}\label{regdesc:EFUSERDRSDATAERR0REG}
    \regfield{EFUSE\_RD\_KEY4\_DATA\_FAIL}{1}{31}{{0}}%
    \regfield{EFUSE\_RD\_KEY4\_DATA\_ERR\_NUM}{3}{28}{{0x0}}%
    \regfield{EFUSE\_RD\_KEY3\_DATA\_FAIL}{1}{27}{{0}}%
    \regfield{EFUSE\_RD\_KEY3\_DATA\_ERR\_NUM}{3}{24}{{0x0}}%
    \regfield{EFUSE\_RD\_KEY2\_DATA\_FAIL}{1}{23}{{0}}%
    \regfield{EFUSE\_RD\_KEY2\_DATA\_ERR\_NUM}{3}{20}{{0x0}}%
    \regfield{EFUSE\_RD\_KEY1\_DATA\_FAIL}{1}{19}{{0}}%
    \regfield{EFUSE\_RD\_KEY1\_DATA\_ERR\_NUM}{3}{16}{{0x0}}%
    \regfield{EFUSE\_RD\_KEY0\_DATA\_FAIL}{1}{15}{{0}}%
    \regfield{EFUSE\_RD\_KEY0\_DATA\_ERR\_NUM}{3}{12}{{0x0}}%
    \regfield{EFUSE\_RD\_USR\_DATA\_FAIL}{1}{11}{{0}}%
    \regfield{EFUSE\_RD\_USR\_DATA\_ERR\_NUM}{3}{8}{{0x0}}%
    \regfield{EFUSE\_RD\_SYS\_PART1\_DATA\_FAIL}{1}{7}{{0}}%
    \regfield{EFUSE\_RD\_SYS\_PART1\_DATA\_ERR\_NUM}{3}{4}{{0x0}}%
    \regfield{EFUSE\_RD\_MAC\_SYS\_FAIL}{1}{3}{{0}}%
    \regfield{EFUSE\_RD\_MAC\_SYS\_ERR\_NUM}{3}{0}{{0x0}}%
    \reglabel{Reset}\regnewline%
    \begin{regdesc}\begin{reglist}
    \label{fielddesc:EFUSERDMACSYSERRNUM}\item [EFUSE\_RD\_MAC\_SYS\_ERR\_NUM] 表示 RD\_MAC\_SYS 中的错误字节数。(RO)
    \label{fielddesc:EFUSERDMACSYSFAIL}\item [EFUSE\_RD\_MAC\_SYS\_FAIL] 表示 RD\_MAC\_SYS 的数据是否烧写错误。\\
    0：无烧写错误，RD\_MAC\_SYS 的数据是可靠的\\
    1：用户数据烧写失败，错误字节数超过 6\\
    (RO)
    \label{fielddesc:EFUSERDSYSPART1DATAERRNUM}\item [EFUSE\_RD\_SYS\_PART1\_DATA\_ERR\_NUM] 表示 RD\_SYS\_PART1\_DATA 中的错误字节数。(RO)
    \label{fielddesc:EFUSERDSYSPART1DATAFAIL}\item [EFUSE\_RD\_SYS\_PART1\_DATA\_FAIL] 表示 RD\_SYS\_PART1\_DATA 的数据是否烧写失败。\\
    0：无烧写错误，RD\_SYS\_PART1\_DATA 的数据是可靠的\\
    1：用户数据烧写失败，错误字节数超过 6\\
    (RO)
    \label{fielddesc:EFUSERDUSRDATAERRNUM}\item [EFUSE\_RD\_USR\_DATA\_ERR\_NUM] 表示 RD\_USR\_DATA 中的错误字节数。 (RO)
    \label{fielddesc:EFUSERDUSRDATAFAIL}\item [EFUSE\_RD\_USR\_DATA\_FAIL] 表示 RD\_USR\_DATA 的数据是否烧写错误。\\
    0：无烧写错误，RD\_USR\_DATA 的数据是可靠的\\
    1：用户数据烧写失败，错误字节数超过 6\\
    (RO)
    \label{fielddesc:EFUSERDKEY0DATAERRNUM}\item [EFUSE\_RD\_KEY0\_DATA\_ERR\_NUM] 表示 RD\_KEY0\_DATA 中的错误字节数。(RO)
    \label{fielddesc:EFUSERDKEY0DATAFAIL}\item [EFUSE\_RD\_KEY0\_DATA\_FAIL] 表示 RD\_KEY0\_DATA 的数据是否烧写错误。\\
    0：无烧写错误，RD\_KEY0\_DATA 的数据是可靠的\\
    1：用户数据烧写失败，错误字节数超过 6\\
    (RO)
    \label{fielddesc:EFUSERDKEY1DATAERRNUM}\item [EFUSE\_RD\_KEY1\_DATA\_ERR\_NUM] 表示 RD\_KEY1\_DATA 中的错误字节数。(RO)
    \label{fielddesc:EFUSERDKEY1DATAFAIL}\item [EFUSE\_RD\_KEY1\_DATA\_FAIL] 表示 RD\_KEY1\_DATA 的数据是否烧写错误。\\
    0：无烧写错误，RD\_KEY1\_DATA 的数据是可靠的\\
    1：用户数据烧写失败，错误字节数超过 6\\
    (RO)
    \label{fielddesc:EFUSERDKEY2DATAERRNUM}\item [EFUSE\_RD\_KEY2\_DATA\_ERR\_NUM] 表示 RD\_KEY2\_DATA 中的错误字节数。(RO)

    \item [见下页……]
    \end{reglist}\end{regdesc}
    \end{register}

    \addtocounter{Regfloat}{-1}
    \begin{register}{H}{EFUSE\_RD\_RS\_DATA\_ERR0\_REG}{0x0190}
    \begin{regdesc}\begin{reglist}
    \item [接上页……]

    \label{fielddesc:EFUSERDKEY2DATAFAIL}\item [EFUSE\_RD\_KEY2\_DATA\_FAIL] 表示 RD\_KEY2\_DATA 的数据是否烧写错误。\\
    0：无烧写错误，RD\_KEY2\_DATA 的数据是可靠的\\
    1：用户数据烧写失败，错误字节数超过 6\\
    (RO)
    \label{fielddesc:EFUSERDKEY3DATAERRNUM}\item [EFUSE\_RD\_KEY3\_DATA\_ERR\_NUM] 表示 RD\_KEY3\_DATA 中的错误字节数。(RO)
    \label{fielddesc:EFUSERDKEY3DATAFAIL}\item [EFUSE\_RD\_KEY3\_DATA\_FAIL] 表示 RD\_KEY3\_DATA 的数据是否烧写错误。\\
    0：无烧写错误，RD\_KEY3\_DATA 的数据是可靠的\\
    1：用户数据烧写失败，错误字节数超过 6\\
    (RO)
    \label{fielddesc:EFUSERDKEY4DATAERRNUM}\item [EFUSE\_RD\_KEY4\_DATA\_ERR\_NUM] 表示 RD\_KEY4\_DATA 中的错误字节数。(RO)
    \label{fielddesc:EFUSERDKEY4DATAFAIL}\item [EFUSE\_RD\_KEY4\_DATA\_FAIL] 表示 RD\_KEY4\_DATA 的数据是否烧写错误。\\
    0：无烧写错误，RD\_KEY4\_DATA 的数据是可靠的\\
    1：用户数据烧写失败，错误字节数超过 6\\
    (RO)
    \end{reglist}\end{regdesc}
    \end{register}


    \begin{register}{H}{EFUSE\_RD\_RS\_DATA\_ERR1\_REG}{0x0194}\label{regdesc:EFUSERDRSDATAERR1REG}
    \regfield{(reserved)}{24}{8}{000000000000000000000000}%
    \regfield{EFUSE\_RD\_SYS\_PART2\_DATA\_FAIL}{1}{7}{{0}}%
    \regfield{EFUSE\_RD\_SYS\_PART2\_DATA\_ERR\_NUM}{3}{4}{{0x0}}%
    \regfield{EFUSE\_RD\_KEY5\_DATA\_FAIL}{1}{3}{{0}}%
    \regfield{EFUSE\_RD\_KEY5\_DATA\_ERR\_NUM}{3}{0}{{0x0}}%
    \reglabel{Reset}\regnewline%
    \begin{regdesc}\begin{reglist}
    \label{fielddesc:EFUSERDKEY5DATAERRNUM}\item [EFUSE\_RD\_KEY5\_DATA\_ERR\_NUM] 表示 RD\_KEY5\_DATA 中的错误字节数。(RO)
    \label{fielddesc:EFUSERDKEY5DATAFAIL}\item [EFUSE\_RD\_KEY5\_DATA\_FAIL] 表示 RD\_KEY5\_DATA 的数据是否烧写错误。\\
    0：无烧写错误，RD\_KEY5\_DATA 的数据是可靠的\\
    1：用户数据烧写失败，错误字节数超过 6\\
    (RO)
    \label{fielddesc:EFUSERDSYSPART2DATAERRNUM}\item [EFUSE\_RD\_SYS\_PART2\_DATA\_ERR\_NUM] 表示 RD\_SYS\_PART2\_DATA 中的错误字节数。(RO)
    \label{fielddesc:EFUSERDSYSPART2DATAFAIL}\item [EFUSE\_RD\_SYS\_PART2\_DATA\_FAIL] 表示 RD\_SYS\_PART2\_DATA 的数据是否烧写错误。\\
    0：无烧写错误，RD\_SYS\_PART2\_DATA 的数据是可靠的\\
    1：用户数据烧写失败，错误字节数超过 6\\
    (RO)
    \end{reglist}\end{regdesc}
    \end{register}


    \begin{register}{H}{EFUSE\_DATE\_REG}{0x0198}\label{regdesc:EFUSEDATEREG}
    \regfield{(reserved)}{4}{28}{0000}%
    \regfield{EFUSE\_DATE}{28}{0}{{0x2401100}}%
    \reglabel{Reset}\regnewline%
    \begin{regdesc}\begin{reglist}
    \label{fielddesc:EFUSEDATE}\item [EFUSE\_DATE] 版本控制寄存器 (R/W)
    \end{reglist}\end{regdesc}
    \end{register}


    \begin{register}{H}{EFUSE\_CLK\_REG}{0x01C8}\label{regdesc:EFUSECLKREG}
    \regfield{(reserved)}{15}{17}{000000000000000}%
    \regfield{EFUSE\_CLK\_EN}{1}{16}{{0}}%
    \regfield{(reserved)}{13}{3}{0000000000000}%
    \regfield{EFUSE\_MEM\_FORCE\_PU}{1}{2}{{0}}%
    \regfield{EFUSE\_MEM\_CLK\_FORCE\_ON}{1}{1}{{1}}%
    \regfield{EFUSE\_MEM\_FORCE\_PD}{1}{0}{{0}}%
    \reglabel{Reset}\regnewline%
    \begin{regdesc}\begin{reglist}
    \label{fielddesc:EFUSEMEMFORCEPD}\item [EFUSE\_MEM\_FORCE\_PD] 配置是否强制 eFuse SRAM 进入节能模式。\\
    1：强制\\
    0：没有作用\\
    (R/W)
    \label{fielddesc:EFUSEMEMCLKFORCEON}\item [EFUSE\_MEM\_CLK\_FORCE\_ON]  配置是否强制激活 eFuse SRAM 时钟信号。\\
    1：强制激活\\
    0：没有作用\\
    (R/W)
    \label{fielddesc:EFUSEMEMFORCEPU}\item [EFUSE\_MEM\_FORCE\_PU] 配置是否强制 eFuse SRAM 进入工作模式。\\
    1：强制\\
    0：没有作用\\
    (R/W)
    \label{fielddesc:EFUSECLKEN}\item [EFUSE\_CLK\_EN] 配置是否强制使能 eFuse 寄存器配置时钟信号。\\
    1：强制使能\\
    0：时钟仅在读写寄存器时开启\\
    (R/W)
    \end{reglist}\end{regdesc}
    \end{register}


    \begin{register}{H}{EFUSE\_CONF\_REG}{0x01CC}\label{regdesc:EFUSECONFREG}
    \regfield{(reserved)}{12}{20}{000000000000}%
    \regfield{EFUSE\_CFG\_ECDSA\_BLK}{4}{16}{{0x0}}%
    \regfield{EFUSE\_OP\_CODE}{16}{0}{{0x00}}%
    \reglabel{Reset}\regnewline%
    \begin{regdesc}\begin{reglist}
    \label{fielddesc:EFUSEOPCODE}\item [EFUSE\_OP\_CODE] 配置操作指令类型。\\
    0x5A5A：烧写操作指令\\
    0x5AA5：读操作指令\\
    其他值：没有作用\\
    (R/W)
    \label{fielddesc:EFUSECFGECDSABLK}\item [EFUSE\_CFG\_ECDSA\_BLK] 配置用于 ECDSA\_DS 密钥输出的区块。(R/W)
    \end{reglist}\end{regdesc}
    \end{register}


    \begin{register}{H}{EFUSE\_DAC\_CONF\_REG}{0x01E8}\label{regdesc:EFUSEDACCONFREG}
    \regfield{(reserved)}{14}{18}{00000000000000}%
    \regfield{EFUSE\_OE\_CLR}{1}{17}{{0}}%
    \regfield{EFUSE\_DAC\_NUM}{8}{9}{{255}}%
    \regfield{EFUSE\_DAC\_CLK\_PAD\_SEL}{1}{8}{{0}}%
    \regfield{EFUSE\_DAC\_CLK\_DIV}{8}{0}{{23}}%
    \reglabel{Reset}\regnewline%
    \begin{regdesc}\begin{reglist}
    \label{fielddesc:EFUSEDACCLKDIV}\item [EFUSE\_DAC\_CLK\_DIV] 配置烧写电压的爬升时钟分频系数。(R/W)
    \label{fielddesc:EFUSEDACCLKPADSEL}\item [EFUSE\_DAC\_CLK\_PAD\_SEL] 保留 (R/W)
    \label{fielddesc:EFUSEDACNUM}\item [EFUSE\_DAC\_NUM] 配置烧写电压上升周期，单位：一个 eFuse 核心时钟周期。(R/W)
    \label{fielddesc:EFUSEOECLR}\item [EFUSE\_OE\_CLR] 降低烧写电压的供电能力。(R/W)
    \end{reglist}\end{regdesc}
    \end{register}


    \begin{register}{H}{EFUSE\_RD\_TIM\_CONF\_REG}{0x01EC}\label{regdesc:EFUSERDTIMCONFREG}
    \regfield{EFUSE\_READ\_INIT\_NUM}{8}{24}{{0x12}}%
    \regfield{EFUSE\_TSUR\_A}{8}{16}{{0x1}}%
    \regfield{EFUSE\_TRD}{8}{8}{{0x2}}%
    \regfield{EFUSE\_THR\_A}{8}{0}{{0x1}}%
    \reglabel{Reset}\regnewline%
    \begin{regdesc}\begin{reglist}
    \label{fielddesc:EFUSETHRA}\item [EFUSE\_THR\_A] 配置读操作的保持时间，单位：一个 eFuse 核心时钟周期。(R/W)
    \label{fielddesc:EFUSETRD}\item [EFUSE\_TRD] 配置读操作的时间，单位：一个 eFuse 核心时钟周期。(R/W)
    \label{fielddesc:EFUSETSURA}\item [EFUSE\_TSUR\_A] 配置读操作的设置时间，单位：一个 eFuse 核心时钟周期。(R/W)
    \label{fielddesc:EFUSEREADINITNUM}\item [EFUSE\_READ\_INIT\_NUM] 配置读取 eFuse 存储器的等待时间，单位：一个 eFuse 核心时钟周期。(R/W)
    \end{reglist}\end{regdesc}
    \end{register}


    \begin{register}{H}{EFUSE\_WR\_TIM\_CONF1\_REG}{0x01F0}\label{regdesc:EFUSEWRTIMCONF1REG}
    \regfield{EFUSE\_THP\_A}{8}{24}{{0x1}}%
    \regfield{EFUSE\_PWR\_ON\_NUM}{16}{8}{{0x3000}}%
    \regfield{EFUSE\_TSUP\_A}{8}{0}{{0x1}}%
    \reglabel{Reset}\regnewline%
    \begin{regdesc}\begin{reglist}
    \label{fielddesc:EFUSETSUPA}\item [EFUSE\_TSUP\_A] 配置烧写的设置时间，单位：一个 eFuse 核心时钟周期。(R/W)
    \label{fielddesc:EFUSEPWRONNUM}\item [EFUSE\_PWR\_ON\_NUM] 配置 VDDQ 的上电时间，单位：一个 eFuse 核心时钟周期。(R/W)
    \label{fielddesc:EFUSETHPA}\item [EFUSE\_THP\_A] 配置烧写的保持时间，单位：一个 eFuse 核心时钟周期。(R/W)
    \end{reglist}\end{regdesc}
    \end{register}


    \begin{register}{H}{EFUSE\_WR\_TIM\_CONF2\_REG}{0x01F4}\label{regdesc:EFUSEWRTIMCONF2REG}
    \regfield{EFUSE\_TPGM}{16}{16}{{0xc8}}%
    \regfield{EFUSE\_PWR\_OFF\_NUM}{16}{0}{{0x190}}%
    \reglabel{Reset}\regnewline%
    \begin{regdesc}\begin{reglist}
    \label{fielddesc:EFUSEPWROFFNUM}\item [EFUSE\_PWR\_OFF\_NUM] 配置 VDDQ 的掉电时间，单位：一个 eFuse 核心时钟周期。(R/W)
    \label{fielddesc:EFUSETPGM}\item [EFUSE\_TPGM] 配置活跃状态的烧写时间，单位：一个 eFuse 核心时钟周期。(R/W)
    \end{reglist}\end{regdesc}
    \end{register}


    \begin{register}{H}{EFUSE\_WR\_TIM\_CONF0\_RS\_BYPASS\_REG}{0x01F8}\label{regdesc:EFUSEWRTIMCONF0RSBYPASSREG}
    \regfield{(reserved)}{11}{21}{00000000000}%
    \regfield{EFUSE\_TPGM\_INACTIVE}{8}{13}{{0x1}}%
    \regfield{EFUSE\_UPDATE}{1}{12}{{0}}%
    \regfield{EFUSE\_BYPASS\_RS\_BLK\_NUM}{11}{1}{{0x0}}%
    \regfield{EFUSE\_BYPASS\_RS\_CORRECTION}{1}{0}{{0}}%
    \reglabel{Reset}\regnewline%
    \begin{regdesc}\begin{reglist}
    \label{fielddesc:EFUSEBYPASSRSCORRECTION}\item [EFUSE\_BYPASS\_RS\_CORRECTION] 配置是否旁路Reed-Solomon 校正。\\
    0：不旁路\\
    1：旁路 \\
    (R/W)
    \label{fielddesc:EFUSEBYPASSRSBLKNUM}\item [EFUSE\_BYPASS\_RS\_BLK\_NUM] 为双重编程操作配置区块编号。 (R/W)
    \label{fielddesc:EFUSEUPDATE}\item [EFUSE\_UPDATE] 配置是否更新多位寄存器信号。\\
    1：更新\\
    0：没有作用\\
    (WT)
    \label{fielddesc:EFUSETPGMINACTIVE}\item [EFUSE\_TPGM\_INACTIVE] 配置非活跃状态的烧写时间，单位：一个 eFuse 核心时钟周期。(R/W)
    \end{reglist}\end{regdesc}
    \end{register}


    \begin{register}{H}{EFUSE\_STATUS\_REG}{0x01D0}\label{regdesc:EFUSESTATUSREG}
    \regfield{(reserved)}{8}{24}{00000000}%
    \regfield{EFUSE\_CUR\_ECDSA\_BLK}{4}{20}{{0x0}}%
    \regfield{EFUSE\_BLK0\_VALID\_BIT\_CNT}{10}{10}{{0x0}}%
    \regfield{EFUSE\_OTP\_VDDQ\_IS\_SW}{1}{9}{{0}}%
    \regfield{EFUSE\_OTP\_PGENB\_SW}{1}{8}{{0}}%
    \regfield{EFUSE\_OTP\_CSB\_SW}{1}{7}{{0}}%
    \regfield{EFUSE\_OTP\_STROBE\_SW}{1}{6}{{0}}%
    \regfield{EFUSE\_OTP\_VDDQ\_C\_SYNC2}{1}{5}{{0}}%
    \regfield{EFUSE\_OTP\_LOAD\_SW}{1}{4}{{0}}%
    \regfield{EFUSE\_STATE}{4}{0}{{0x0}}%
    \reglabel{Reset}\regnewline%
    \begin{regdesc}\begin{reglist}
    \label{fielddesc:EFUSESTATE}\item [EFUSE\_STATE] 表示 eFuse 状态机的状态。(RO)
    \label{fielddesc:EFUSEOTPLOADSW}\item [EFUSE\_OTP\_LOAD\_SW] 表示 OTP\_LOAD\_SW 的值。(RO)
    \label{fielddesc:EFUSEOTPVDDQCSYNC2}\item [EFUSE\_OTP\_VDDQ\_C\_SYNC2] 表示 OTP\_VDDQ\_C\_SYNC2 的值。(RO)
    \label{fielddesc:EFUSEOTPSTROBESW}\item [EFUSE\_OTP\_STROBE\_SW] 表示 OTP\_STROBE\_SW 的值。(RO)
    \label{fielddesc:EFUSEOTPCSBSW}\item [EFUSE\_OTP\_CSB\_SW] 表示 OTP\_CSB\_SW 的值。(RO)
    \label{fielddesc:EFUSEOTPPGENBSW}\item [EFUSE\_OTP\_PGENB\_SW] 表示 OTP\_PGENB\_SW 的值。(RO)
    \label{fielddesc:EFUSEOTPVDDQISSW}\item [EFUSE\_OTP\_VDDQ\_IS\_SW] 表示 OTP\_VDDQ\_IS\_SW 的值。(RO)
    \label{fielddesc:EFUSEBLK0VALIDBITCNT}\item [EFUSE\_BLK0\_VALID\_BIT\_CNT] 表示有效块位数。(RO)
    \label{fielddesc:EFUSECURECDSABLK}\item [EFUSE\_CUR\_ECDSA\_BLK] 表示用于 ECDSA\_DS 密钥输出的区块。(RO)
    \end{reglist}\end{regdesc}
    \end{register}


    \begin{register}{H}{EFUSE\_CMD\_REG}{0x01D4}\label{regdesc:EFUSECMDREG}
    \regfield{(reserved)}{26}{6}{00000000000000000000000000}%
    \regfield{EFUSE\_BLK\_NUM}{4}{2}{{0x0}}%
    \regfield{EFUSE\_PGM\_CMD}{1}{1}{{0}}%
    \regfield{EFUSE\_READ\_CMD}{1}{0}{{0}}%
    \reglabel{Reset}\regnewline%
    \begin{regdesc}\begin{reglist}
    \label{fielddesc:EFUSEREADCMD}\item [EFUSE\_READ\_CMD] 配置是否发送读指令。\\
    1：发送\\
    0：没有作用\\
    (R/W/SC)
    \label{fielddesc:EFUSEPGMCMD}\item [EFUSE\_PGM\_CMD]  配置是否发送烧写指令。\\
    1：发送\\
    0：没有作用\\
    (R/W/SC)
    \label{fielddesc:EFUSEBLKNUM}\item [EFUSE\_BLK\_NUM] 表示待烧写块的序号。值 0 \verb+~+ 10 分别对应 BLOCK0 \verb+~+ BLOCK10。(R/W)
    \end{reglist}\end{regdesc}
    \end{register}


    \begin{register}{H}{EFUSE\_INT\_RAW\_REG}{0x01D8}\label{regdesc:EFUSEINTRAWREG}
    \regfield{(reserved)}{30}{2}{000000000000000000000000000000}%
    \regfield{EFUSE\_PGM\_DONE\_INT\_RAW}{1}{1}{{0}}%
    \regfield{EFUSE\_READ\_DONE\_INT\_RAW}{1}{0}{{0}}%
    \reglabel{Reset}\regnewline%
    \begin{regdesc}\begin{reglist}
    \label{fielddesc:EFUSEREADDONEINTRAW}\item [EFUSE\_READ\_DONE\_INT\_RAW] 读取完成的原始中断状态。(R/SS/WTC)
    \label{fielddesc:EFUSEPGMDONEINTRAW}\item [EFUSE\_PGM\_DONE\_INT\_RAW] 烧写完成的原始中断状态。(R/SS/WTC)
    \end{reglist}\end{regdesc}
    \end{register}


    \begin{register}{H}{EFUSE\_INT\_ST\_REG}{0x01DC}\label{regdesc:EFUSEINTSTREG}
    \regfield{(reserved)}{30}{2}{000000000000000000000000000000}%
    \regfield{EFUSE\_PGM\_DONE\_INT\_ST}{1}{1}{{0}}%
    \regfield{EFUSE\_READ\_DONE\_INT\_ST}{1}{0}{{0}}%
    \reglabel{Reset}\regnewline%
    \begin{regdesc}\begin{reglist}
    \label{fielddesc:EFUSEREADDONEINTST}\item [EFUSE\_READ\_DONE\_INT\_ST] 读取完成的屏蔽中断状态。(RO)
    \label{fielddesc:EFUSEPGMDONEINTST}\item [EFUSE\_PGM\_DONE\_INT\_ST] 烧写完成的屏蔽中断状态。(RO)
    \end{reglist}\end{regdesc}
    \end{register}


    \begin{register}{H}{EFUSE\_INT\_ENA\_REG}{0x01E0}\label{regdesc:EFUSEINTENAREG}
    \regfield{(reserved)}{30}{2}{000000000000000000000000000000}%
    \regfield{EFUSE\_PGM\_DONE\_INT\_ENA}{1}{1}{{0}}%
    \regfield{EFUSE\_READ\_DONE\_INT\_ENA}{1}{0}{{0}}%
    \reglabel{Reset}\regnewline%
    \begin{regdesc}\begin{reglist}
    \label{fielddesc:EFUSEREADDONEINTENA}\item [EFUSE\_READ\_DONE\_INT\_ENA] 写 1 使能读取完成的中断。(R/W)
    \label{fielddesc:EFUSEPGMDONEINTENA}\item [EFUSE\_PGM\_DONE\_INT\_ENA] 写 1 使能烧写完成的中断。(R/W)
    \end{reglist}\end{regdesc}
    \end{register}


    \begin{register}{H}{EFUSE\_INT\_CLR\_REG}{0x01E4}\label{regdesc:EFUSEINTCLRREG}
    \regfield{(reserved)}{30}{2}{000000000000000000000000000000}%
    \regfield{EFUSE\_PGM\_DONE\_INT\_CLR}{1}{1}{{0}}%
    \regfield{EFUSE\_READ\_DONE\_INT\_CLR}{1}{0}{{0}}%
    \reglabel{Reset}\regnewline%
    \begin{regdesc}\begin{reglist}
    \label{fielddesc:EFUSEREADDONEINTCLR}\item [EFUSE\_READ\_DONE\_INT\_CLR] 写 1 清除读取完成的中断。(WT)
    \label{fielddesc:EFUSEPGMDONEINTCLR}\item [EFUSE\_PGM\_DONE\_INT\_CLR] 写 1 清除烧写完成的中断。(WT)
    \end{reglist}\end{regdesc}
    \end{register}
